\documentclass[12pt, letterpaper]{article}

% ── Paquetes ──────────────────────────────────────────────────────────────────
\usepackage[utf8]{inputenc}
\usepackage[spanish]{babel}
\usepackage[margin=2.5cm]{geometry}
\usepackage{graphicx}
\usepackage{xcolor}
\usepackage{booktabs}
\usepackage{tabularx}
\usepackage{enumitem}
\usepackage{titlesec}
\usepackage{fancyhdr}
\usepackage{mdframed}
\usepackage{fontawesome5}
\usepackage{hyperref}
\usepackage{draftwatermark}

% ── Marca de agua BORRADOR ────────────────────────────────────────────────────
\SetWatermarkText{BORRADOR}
\SetWatermarkScale{3.5}
\SetWatermarkColor[gray]{0.88}

% ── Colores institucionales UASD ─────────────────────────────────────────────
\definecolor{azuluasd}{RGB}{0, 63, 135}
\definecolor{doradouasd}{RGB}{253, 185, 19}
\definecolor{grisclaro}{RGB}{244, 246, 251}
\definecolor{verdexito}{RGB}{21, 87, 36}
\definecolor{verdebg}{RGB}{212, 237, 218}
\definecolor{amarillobg}{RGB}{255, 243, 205}
\definecolor{rojoalerta}{RGB}{114, 28, 36}

% ── Hipervínculos ─────────────────────────────────────────────────────────────
\hypersetup{
    colorlinks=true,
    linkcolor=azuluasd,
    urlcolor=azuluasd,
}

% ── Encabezado y pie de página ────────────────────────────────────────────────
\pagestyle{fancy}
\fancyhf{}
\fancyhead[L]{\textcolor{azuluasd}{\textbf{CONCALAB-UASD}}}
\fancyhead[R]{\textcolor{gray}{\small Control de Calidad de Laboratorios}}
\fancyfoot[C]{\textcolor{gray}{\small\thepage}}
\fancyfoot[R]{\textcolor{gray}{\small\textit{Documento interno — BORRADOR 2026}}}
\renewcommand{\headrulewidth}{2pt}
\renewcommand{\headrule}{\hbox to\headwidth{\color{azuluasd}\leaders\hrule height \headrulewidth\hfill}}

% ── Títulos de sección ────────────────────────────────────────────────────────
\titleformat{\section}
  {\large\bfseries\color{azuluasd}}
  {\thesection.}{0.6em}{}
  [\color{doradouasd}\titlerule[1.5pt]]

\titleformat{\subsection}
  {\normalsize\bfseries\color{azuluasd}}
  {\thesubsection.}{0.5em}{}

% ── Cajas de aviso ────────────────────────────────────────────────────────────
\newmdenv[
  backgroundcolor=amarillobg,
  linecolor=doradouasd,
  linewidth=1.5pt,
  leftline=true, rightline=false, topline=false, bottomline=false,
  innerleftmargin=10pt, innerrightmargin=8pt,
  innertopmargin=6pt, innerbottommargin=6pt,
  skipabove=8pt, skipbelow=8pt,
]{cajaaviso}

\newmdenv[
  backgroundcolor=verdebg,
  linecolor=verdexito,
  linewidth=1.5pt,
  leftline=true, rightline=false, topline=false, bottomline=false,
  innerleftmargin=10pt, innerrightmargin=8pt,
  innertopmargin=6pt, innerbottommargin=6pt,
  skipabove=8pt, skipbelow=8pt,
]{cajaexito}

\newmdenv[
  backgroundcolor=grisclaro,
  linecolor=azuluasd,
  linewidth=1.5pt,
  leftline=true, rightline=false, topline=false, bottomline=false,
  innerleftmargin=10pt, innerrightmargin=8pt,
  innertopmargin=6pt, innerbottommargin=6pt,
  skipabove=8pt, skipbelow=8pt,
]{cajainfo}

% ─────────────────────────────────────────────────────────────────────────────
\begin{document}

% ── Portada ───────────────────────────────────────────────────────────────────
\begin{titlepage}
  \centering
  \vspace*{1.5cm}

  {\Huge\bfseries\color{azuluasd} CONCALAB-UASD}\\[0.4cm]
  {\large\color{gray} Universidad Autónoma de Santo Domingo}\\[0.2cm]
  {\large\color{gray} Control de Calidad de Laboratorios}

  \vspace{1.5cm}
  \noindent\color{doradouasd}\rule{\linewidth}{2pt}
  \vspace{0.5cm}

  {\LARGE\bfseries\color{azuluasd}
    Instructivo para el Reporte de Resultados\\[0.3cm]
    de Ensayos de Aptitud}

  \vspace{0.5cm}
  \noindent\color{doradouasd}\rule{\linewidth}{2pt}

  \vspace{1.2cm}
  \colorbox{doradouasd}{%
    \parbox{6cm}{\centering\large\bfseries\color{azuluasd} BORRADOR}%
  }

  \vfill

  \begin{tabular}{ll}
    \textbf{Ronda:}    & Primera Ronda de Ensayo de Aptitud 2026 (EA-001-2026) \\[4pt]
    \textbf{Versión:}  & 1.0 — Borrador \\[4pt]
    \textbf{Fecha:}    & Mayo 2026 \\[4pt]
    \textbf{Elaborado por:} & CONCALAB-UASD \\
  \end{tabular}

  \vspace{1.5cm}
  {\small\color{gray} Este documento es de uso interno y no debe distribuirse sin autorización.}
\end{titlepage}

% ── Cuerpo ────────────────────────────────────────────────────────────────────
\tableofcontents
\newpage

% ─────────────────────────────────────────────────────────────────────────────
\section{Requisitos previos}

Antes de iniciar el proceso de reporte, asegúrese de contar con lo siguiente:

\begin{itemize}[leftmargin=1.5em]
  \item Credenciales de acceso (correo electrónico y contraseña) asignadas por CONCALAB-UASD.
  \item Resultados de la ronda activa disponibles y verificados.
  \item Conexión a internet estable.
  \item Navegador actualizado: Google Chrome, Mozilla Firefox o Microsoft Edge.
\end{itemize}

% ─────────────────────────────────────────────────────────────────────────────
\section{Ingresar al formulario}

\begin{enumerate}[leftmargin=1.5em]
  \item Abra su navegador.
  \item Ingrese a la página web oficial de CONCALAB-UASD.
  \item En el menú de navegación, haga clic en \textbf{Servicios → Reportar Resultados}.
\end{enumerate}

% ─────────────────────────────────────────────────────────────────────────────
\section{Iniciar sesión}

\begin{enumerate}[leftmargin=1.5em]
  \item Ingrese su \textbf{correo electrónico} y \textbf{contraseña} asignados por CONCALAB-UASD.
  \item Haga clic en \textbf{Iniciar Sesión}.
\end{enumerate}

\begin{cajaaviso}
\textbf{⚠ Importante:} Si no recuerda su contraseña, contacte a CONCALAB-UASD antes de intentar acceder varias veces. Después de múltiples intentos fallidos, el sistema bloqueará el acceso temporalmente.
\end{cajaaviso}

% ─────────────────────────────────────────────────────────────────────────────
\section{Verificar la información pre-cargada}

Una vez autenticado, el formulario mostrará automáticamente:

\begin{itemize}[leftmargin=1.5em]
  \item \textbf{Nombre de su laboratorio} y código asignado (ej: \texttt{AG4}).
  \item \textbf{Código de la ronda activa} (ej: \texttt{EA-001-2026}).
\end{itemize}

Estos campos \textbf{no son editables}. Si detecta algún error en esta información, comuníquese con CONCALAB-UASD antes de continuar.

% ─────────────────────────────────────────────────────────────────────────────
\section{Completar la información general}

Llene los siguientes campos:

\vspace{0.5em}
\begin{tabularx}{\linewidth}{>{\bfseries}lX}
  \toprule
  Campo & Descripción \\
  \midrule
  Fecha del Reporte    & Fecha en que reporta los resultados (formato DD/MM/AAAA). \\[4pt]
  Correo Electrónico   & Se completa automáticamente. Puede modificarlo si es necesario. \\
  \bottomrule
\end{tabularx}

% ─────────────────────────────────────────────────────────────────────────────
\section{Ingresar resultados de Química Clínica}

La tabla muestra los analitos de la ronda activa. Para cada analito que haya procesado, complete las siguientes columnas:

\vspace{0.5em}
\begin{tabularx}{\linewidth}{>{\bfseries}lX}
  \toprule
  Columna & Descripción \\
  \midrule
  Instrumento & Nombre del equipo utilizado (ej: Cobas 6000, Mindray BS-430). \\[4pt]
  Método      & Método analítico empleado (ej: Enzimático, Colorimétrico). \\[4pt]
  Resultado   & Valor numérico obtenido. \\[4pt]
  Unidad      & Unidad de medida (ej: mg/dL, U/L, mmol/L). \\
  \bottomrule
\end{tabularx}

\begin{cajainfo}
\textbf{Consejo:} Use el botón \textbf{⬇ Copiar} en el encabezado de la columna Instrumento para repetir automáticamente el mismo equipo en todos los analitos.
\end{cajainfo}

% ─────────────────────────────────────────────────────────────────────────────
\section{Ingresar resultados de Uroanálisis}

Complete la tabla de Uroanálisis de la misma forma que la de Química Clínica. Los resultados pueden ser cualitativos (ej: Negativo, +, ++, +++).

% ─────────────────────────────────────────────────────────────────────────────
\section{Comentarios adicionales (opcional)}

Si desea agregar alguna observación sobre los resultados o las condiciones del análisis, utilice el campo \textbf{Comentarios o Notas Adicionales}.

% ─────────────────────────────────────────────────────────────────────────────
\section{Enviar los resultados}

\begin{enumerate}[leftmargin=1.5em]
  \item Haga clic en el botón \textbf{Enviar Resultados}.
  \item Aparecerá una ventana de \textbf{confirmación} con el resumen de los datos:
    \begin{itemize}
      \item Nombre del laboratorio, código y ronda.
      \item Fecha y correo de contacto.
      \item Listado completo de analitos con sus resultados.
    \end{itemize}
  \item Revise cuidadosamente el resumen:
    \begin{itemize}
      \item Si todo está correcto → haga clic en \textbf{Confirmar y Enviar}.
      \item Si necesita corregir algo → haga clic en \textbf{Revisar de nuevo}.
    \end{itemize}
\end{enumerate}

\begin{cajaaviso}
\textbf{⚠ Importante:} Una vez confirmado el envío, los resultados \textbf{no podrán ser modificados} desde el portal. Cualquier corrección debe solicitarse directamente a CONCALAB-UASD.
\end{cajaaviso}

% ─────────────────────────────────────────────────────────────────────────────
\section{Confirmación del envío}

Si el envío fue exitoso, verá el siguiente mensaje en pantalla:

\begin{cajaexito}
\textbf{✓ ¡Reporte Enviado Exitosamente!}\\
Los datos han sido registrados en nuestra base de datos segura.\\
Se ha enviado una copia al correo electrónico registrado.
\end{cajaexito}

Además, recibirá un \textbf{correo electrónico} con el resumen completo de los resultados reportados como constancia.

% ─────────────────────────────────────────────────────────────────────────────
\section{Situaciones especiales}

\subsection{El formulario no está disponible}
Si aparece el mensaje \textit{``El formulario de reporte no está habilitado en este momento''}, la ronda aún no ha sido abierta por CONCALAB-UASD. Espere la notificación oficial.

\subsection{Ya reportó en esta ronda}
Si aparece el mensaje \textit{``Ya enviaste resultados para EA-XXX-XXXX''}, su reporte fue registrado previamente. Si necesita una corrección, contacte a CONCALAB-UASD.

\subsection{Problemas de acceso}
Si su cuenta no es reconocida o aparece un mensaje de error, comuníquese con CONCALAB-UASD para verificar su registro en el sistema.

% ─────────────────────────────────────────────────────────────────────────────
\section{Contacto}

\begin{cajainfo}
\textbf{CONCALAB-UASD — Control de Calidad de Laboratorios}\\
Universidad Autónoma de Santo Domingo (UASD)\\
Correo: \href{mailto:concalab@uasd.edu.do}{concalab@uasd.edu.do}
\end{cajainfo}

\vfill
{\small\color{gray}\centering Documento interno — CONCALAB-UASD — BORRADOR 2026. No distribuir sin autorización.}

\end{document}
